\section{根理想}\label{sec:radical-ideals}

上一节说明了 $K[x_{1},\ldots,x_{n}]$ 中并非每个理想都能作为 $\mathbb{A}^{n}$ 中某个子集的消没理想. 特别地, 在 Example~\ref{ex:square-ideal-not-vanishing-ideal} 中, 我们看到理想 $\langle x^{2}\rangle\subseteq K[x]$ 不是消没理想. 那么, 怎样判断给定的理想是否为某个 $X$ 的 $\mathcal{I}(X)$ 呢? 本节来考察消没理想的一个重要性质: 消没理想都是根理想.

\begin{Definition}{根理想}\leavevmode\label{def:radical-ideal}
  若对所有 $a\in R$ 都有
  \[
    a^{m}\in I\text{ 对某个整数 }m>0\text{ 成立}
    \quad\Longrightarrow\quad a\in I,
  \]
  则称理想 $I\subseteq R$ 为\textbf{根理想}.
\end{Definition}

换句话说, 一个理想是根理想, 当且仅当它对开方封闭: 只要一个元素的某次幂属于理想, 该元素本身就必须属于理想. 注意, Example~\ref{ex:square-ideal-not-vanishing-ideal} 中的理想 $\langle x^{2}\rangle\subseteq K[x]$ 不是根理想, 因为 $x^{2}\in\langle x^{2}\rangle$, 但 $x\notin\langle x^{2}\rangle$. 下面的结果表明, 所有消没理想都具有这个性质.

\begin{Proposition}{消没理想都是根理想}\leavevmode\label{prop:vanishing-ideals-are-radical}
  对任意集合 $X\subseteq\mathbb{A}^{n}$, 消没理想 $\mathcal{I}(X)$ 是根理想.
\end{Proposition}

特别地, 由 Proposition~\ref{prop:vanishing-ideals-are-radical} 和 $\langle x^{2}\rangle$ 不是根理想这一事实可知, $\langle x^{2}\rangle$ 不是消没理想, 正如我们在 Example~\ref{ex:square-ideal-not-vanishing-ideal} 中看到的那样.

\begin{Proof}\unskip\textbf{ of Proposition~\ref{prop:vanishing-ideals-are-radical}}\quad
  设 $X\subseteq\mathbb{A}^{n}$ 是一个子集, 并假设存在正整数 $m$ 使得 $f^{m}\in\mathcal{I}(X)$, 换句话说, 对所有 $a\in X$ 都有 $f^{m}(a)=0$. 那么
  \[
    0=f^{m}(a)=\bigl(f(a)\bigr)^{m}\in K.
  \]
  由于 $K$ 是域, 它没有零因子. 特别地, $\bigl(f(a)\bigr)^{m}=0$ 当且仅当 $f(a)=0$. 因而对所有 $a\in X$, 都有 $f(a)=0$, 所以 $f\in\mathcal{I}(X)$.
\end{Proof}

特别地, Proposition~\ref{prop:vanishing-ideals-are-radical} 表明根理想在代数几何中起核心作用. 因此, 本节余下的篇幅将集中介绍与根理想有关的一些基本概念.

根据定义, 若一个理想对开方不封闭, 它就不是根理想. 因此, 你或许会想到, 只要补上缺少的根, 就能构造一个根理想. 下面的定义将这一构造精确化.

\begin{Definition}{理想的根}\leavevmode\label{def:radical-of-an-ideal}
  设 $I\subseteq R$ 是一个理想. $I$ 的\textbf{根}为
  \[
    \sqrt{I}=\{a\in R\mid a^{m}\in I\text{ 对某个 }m>0\text{ 成立}\}.
  \]
\end{Definition}

如下一个结果所示, 这个\textquotedbl{}补上缺少的根\textquotedbl{}的过程确实给出一个根理想.

\begin{Proposition}{理想的根是根理想}\leavevmode\label{prop:radicals-are-radical}
  若 $I\subseteq R$ 是一个理想, 则 $\sqrt{I}$ 是根理想.
\end{Proposition}

\begin{Proof}
  $\sqrt{I}$ 是一个理想的证明见 Exercise~\ref{ex:radical-is-an-ideal}. 为了证明 $\sqrt{I}$ 是根理想, 假设对某个整数 $k>0$, 有 $a^{k}\in\sqrt{I}$, 需要证明 $a\in\sqrt{I}$. 根据 $\sqrt{I}$ 的定义, 有
  \[
    a^{k}\in\sqrt{I}\quad\Longrightarrow\quad(a^{k})^{m}\in I
      \text{ 对某个整数 }m>0\text{ 成立}.
  \]
  换句话说, $a^{(km)}\in I$, 因而 $a\in\sqrt{I}$.
\end{Proof}

如果 $I$ 是根理想, 那么它已经对开方封闭, 所以 $I=\sqrt{I}$. 反过来, 若 $I=\sqrt{I}$, 则 Proposition~\ref{prop:radicals-are-radical} 蕴含 $I$ 是根理想. 这样我们证明了以下对根理想的有用刻画.

\begin{Corollary}{根理想的刻画}\leavevmode\label{cor:characterize-radical-ideals}
  理想 $I\subseteq R$ 是根理想, 当且仅当 $I=\sqrt{I}$.
\end{Corollary}

下面举几个例子, 帮助读者熟悉根理想.

\begin{Example}{}\leavevmode\label{ex:radical-of-square-variable}
  若 $R=K[x]$ 且 $I=\langle x^{2}\rangle$, 我们已经看到 $I$ 不是根理想. 若要通过加入 $I$ 中元素的所有可能的根来扩大 $I$, 应当加些什么呢? 下面证明
  \[
    \sqrt{I}=\langle x\rangle.
  \]
  首先, 由于 $x^{2}\in I$, 故 $x\in\sqrt{I}$. 因为 $\sqrt{I}$ 是理想, 由 $x\in\sqrt{I}$ 可知 $\langle x\rangle\subseteq\sqrt{I}$, 这就证明了一个方向的包含关系.

  反过来, 假设 $f\notin\langle x\rangle$. 这意味着 $f$ 的常数项非零. 因而对所有整数 $m>0$, $f^{m}$ 的常数项都非零, 所以 $f^{m}\notin I$. 因此 $f\notin\sqrt{I}$, 这就用逆否命题证明了 $\sqrt{I}\subseteq\langle x\rangle$.
\end{Example}

\begin{Example}{}\leavevmode\label{ex:radical-of-twelve}
  若 $R=\mathbb{Z}$ 且 $I=\langle12\rangle$, 那么 $I$ 不是根理想, 因为 $6^{2}\in I$, 但 $6\notin I$. 事实上,
  \[
    \sqrt{I}=\langle6\rangle.
  \]
  为证 $\sqrt{I}\subseteq\langle6\rangle$, 取 $r\in\sqrt{I}$. 则对某个正整数 $m$ 和某个整数 $k$, 有 $r^{m}=12k$. 由于 $2$ 是素数且 $2\mid r^{m}$, 根据 Euclid 引理可得 $2\mid r$. 同理可得 $3\mid r$. 因为 $r$ 同时被 $2$ 和 $3$ 整除, 所以 $r$ 被 $6$ 整除, 即 $r\in\langle6\rangle$, 这就证明了一个方向的包含关系.

  反过来, 若 $r\in\langle6\rangle$, 则对某个整数 $k$, 有 $r=6k$. 因此
  \[
    r^{2}=36k^{2}=12(3k^{2}),
  \]
  所以 $r^{2}\in I$. 由此得 $r\in\sqrt{I}$, 验证了 $\langle6\rangle\subseteq\sqrt{I}$.
\end{Example}

比较上一例中的根理想 $\langle6\rangle$ 和非根理想 $\langle12\rangle$, 可见 $6$ 是无平方因子的, 换句话说, 它不被任何素数的平方整除. 而 $12$ 不是无平方因子的, 它被 $2^{2}$ 整除. 与 Example~\ref{ex:radical-of-twelve} 类似的论证表明, $\langle n\rangle\subseteq\mathbb{Z}$ 是根理想当且仅当 $n$ 无平方因子.

在代数几何的背景下, 能够判断理想 $I\subseteq K[x_{1},\ldots,x_{n}]$ 何时是根理想会很有用. 对一般的理想而言, 这可能相当困难. 不过, 当 $I$ 是主理想时, 情形与整数的情形类似. 为了准确叙述结果, 取非常数多项式 $f\in K[x_{1},\ldots,x_{n}]$, 并考虑其一个不可约分解:
\[
  f=p_{1}\cdots p_{m}.
\]
由不可约性可知, $p_{i}\mid p_{j}$ 当且仅当两者相差一个常数因子. 把相差常数因子的各项归在一起, 可以写成
\begin{equation}\label{eq:distinct-irreducible-factorization}
  f=aq_{1}^{k_{1}}\cdots q_{\ell}^{k_{\ell}}
\end{equation}
其中 $a\in K$ 非零, 每个 $q_{i}$ 都不可约, 且 $i\ne j$ 时 $q_{i}\nmid q_{j}$. 我们称 \eqref{eq:distinct-irreducible-factorization} 为 $f$ 的\textbf{不同不可约因子分解}, 称各 $q_{i}$ 为 $f$ 的\textbf{不同不可约因子}. 由分解的唯一性可知, 除去重新排列和乘以常数的差别, 不同不可约因子是唯一的.

下面的结果用这些因子来描述主理想的根.

\begin{Proposition}{主理想的根}\leavevmode\label{prop:radicals-of-principal-ideals}
  若 $f\in K[x_{1},\ldots,x_{n}]$ 是非常数多项式, 不同不可约因子为 $q_{1},\ldots,q_{\ell}$, 则
  \[
    \sqrt{\langle f\rangle}=\langle q_{1}\cdots q_{\ell}\rangle.
  \]
  特别地, $\langle f\rangle$ 是根理想, 当且仅当 $f$ 不被任何非常数多项式的平方整除.
\end{Proposition}

\begin{Proof}
  我们证明等式两边互相包含, 将充要条件的推导留给 Exercise~\ref{ex:principal-radical-iff-square-free}.

  ($\subseteq$) 假设 $g\in\sqrt{\langle f\rangle}$. 那么对某个 $m>0$ 和某个 $h\in K[x_{1},\ldots,x_{n}]$, 有 $g^{m}=hf$. 特别地, 这意味着对每个 $i$, 都有 $q_{i}\mid g^{m}$. 由不同不可约因子的唯一性可知, 每个 $q_{i}$ 必定是 $g^{m}$ 的不同不可约因子之一. 因为 $g^{m}$ 的不同不可约因子与 $g$ 的相同, 所以每个 $q_{i}$ 也必定是 $g$ 的不同不可约因子之一. 从而 $q_{1}\cdots q_{\ell}\mid g$, 所以 $g\in\langle q_{1}\cdots q_{\ell}\rangle$.

  ($\supseteq$) 假设 $g\in\langle q_{1}\cdots q_{\ell}\rangle$, 写成 $g=hq_{1}\cdots q_{\ell}$, 其中 $h$ 为某个多项式. 对于不同不可约因子分解
  \[
    f=aq_{1}^{k_{1}}\cdots q_{\ell}^{k_{\ell}},
  \]
  令 $m=\max\{k_{1},\ldots,k_{\ell}\}$. 于是
  \[
    g^{m}=h^{m}q_{1}^{m}\cdots q_{\ell}^{m}
      =\bigl(h^{m}q_{1}^{m-k_{1}}\cdots q_{\ell}^{m-k_{\ell}}a^{-1}\bigr)f
      \in\langle f\rangle,
  \]
  因而 $g\in\sqrt{\langle f\rangle}$.
\end{Proof}

\begin{Example}{{\textmd{$\langle x^{2}+y^{2}\rangle\subseteq\mathbb{C}[x,y]$} 是根理想}}\leavevmode\label{ex:sum-of-squares-complex-radical}
  $x^{2}+y^{2}$ 在 $\mathbb{C}[x,y]$ 中的不可约分解为
  \[
    x^{2}+y^{2}=(x-\mathrm{i}y)(x+\mathrm{i}y).
  \]
  因为两个不可约因子不同, 所以 $x^{2}+y^{2}$ 无平方因子. 因此, 理想 $\langle x^{2}+y^{2}\rangle\subseteq\mathbb{C}[x,y]$ 是根理想.
\end{Example}

\begin{Example}{不是消没理想的根理想}\leavevmode\label{ex:real-radical-not-vanishing}
  在实数上, 多项式 $x^{2}+y^{2}\in\mathbb{R}[x,y]$ 不可约. 因此由 Proposition~\ref{prop:radicals-of-principal-ideals}, 理想 $\langle x^{2}+y^{2}\rangle$ 是根理想. 然而, 由于 $x^{2}+y^{2}$ 仅在原点为零, 可得
  \[
    \mathcal{I}\bigl(\mathcal{V}(\langle x^{2}+y^{2}\rangle)\bigr)
     =\mathcal{I}(\{(0,0)\})=\langle x,y\rangle\supsetneq\langle x^{2}+y^{2}\rangle.
  \]
  因此 Proposition~\ref{prop:compose-vanishing-operators} 蕴含 $\langle x^{2}+y^{2}\rangle$ 不是消没理想. 特别地, 这个例子说明 Proposition~\ref{prop:vanishing-ideals-are-radical} 的逆命题在 $\mathbb{R}$ 上不成立: 根理想不一定是消没理想.
\end{Example}

在研究环的过程中, 我们现在已经遇到三类特殊理想: 极大理想, 素理想和根理想. 我们已经知道极大理想是素理想. 下面的结果把根理想也纳入这个层次中.

\begin{Proposition}{素理想和极大理想都是根理想}\leavevmode\label{prop:prime-ideals-are-radical}
  每个素理想都是根理想.
\end{Proposition}

\begin{Proof}
  为证明逆否命题, 假设 $I$ 是一个不是根理想的理想. 选取 $a\notin I$, 使对某个 $m>1$ 有 $a^{m}\in I$. 若 $m_{0}$ 是使 $a^{m_{0}}\in I$ 成立的大于 $1$ 的最小整数, 则 $a$ 和 $a^{m_{0}-1}$ 都不属于 $I$, 但它们的乘积 $a^{m_{0}}$ 属于 $I$. 因此 $I$ 不是素理想.
\end{Proof}

概括起来, 对任意环 $R$, 有如下理想的层次:
\[
  \{\text{理想}\}\supseteq\{\text{根理想}\}\supseteq
  \{\text{素理想}\}\supseteq\{\text{极大理想}\}.
\]
这些包含关系并不对每个环都是严格的, 但对域上的多元多项式环, 它们是严格的 (Exercise~\ref{ex:strict-hierarchy-of-ideals}).

本节最后, 回到我们的总目标: 在代数对象和几何对象之间建立双射. 由于对于所有仿射簇 $X\subseteq\mathbb{A}^{n}$, 都有 $\mathcal{V}(\mathcal{I}(X))=X$ (Proposition~\ref{prop:compose-vanishing-operators}), 而且 $\mathcal{I}(X)$ 是根理想 (Proposition~\ref{prop:vanishing-ideals-are-radical}), 所以限制定义域后 $\mathcal{V}$ 算子仍然是满射:
\begin{equation}\label{eq:radical-ideal-vanishing-set-surjection}
  \mathcal{V}:\{K[x_{1},\ldots,x_{n}]\text{ 中的根理想}\}
   \longrightarrow\{\mathbb{A}^{n}\text{ 中的仿射簇}\}.
\end{equation}
我们或许会乐观地希望 \eqref{eq:radical-ideal-vanishing-set-surjection} 就是要寻找的代数对象与几何对象之间的双射. 遗憾的是, Example~\ref{ex:real-radical-not-vanishing} 给出了一个反例: $\langle x^{2}+y^{2}\rangle$ 与 $\langle x,y\rangle$ 是 $\mathbb{R}[x,y]$ 中不同的根理想, 却有同一个零点集 $\{(0,0)\}\in\mathbb{A}^{2}_{\mathbb{R}}$. 不过, 如果再对基域 $K$ 加一个假设, 即 $K$ 是代数闭域, 那么 \eqref{eq:radical-ideal-vanishing-set-surjection} 的确就是我们需要的双射. 这个结果是零点定理的推论; 下一节转而讨论这个定理.

\BookExerciseHeading
\begin{BookExercises}
  \item\label{ex:compute-sample-radicals} 判断下列理想哪些是根理想. 对不是根理想的, 求它们的根.
    \begin{enumerate}[label=(\alph*)]
      \item $\langle4\rangle\subseteq\mathbb{Z}$
      \item $\langle6\rangle\subseteq\mathbb{Z}$
      \item $\langle18\rangle\subseteq\mathbb{Z}$
      \item $\langle x^{2}+y^{2}-1\rangle\subseteq\mathbb{R}[x,y]$
      \item $\langle x^{2},y^{3}\rangle\subseteq\mathbb{R}[x,y]$
      \item $\langle y-x^{2},y\rangle\subseteq\mathbb{R}[x,y]$
      \item $\langle x^{2}-y^{2}\rangle\subseteq\mathbb{R}[x,y]$
    \end{enumerate}
  \item\label{ex:radical-is-an-ideal} 设 $I\subseteq R$ 是一个理想. 证明 $\sqrt{I}\subseteq R$ 是一个理想. (提示: 用二项式定理证明 $\sqrt{I}$ 对加法和减法封闭.)
  \item\label{ex:principal-radical-iff-square-free} 证明主理想 $\langle f\rangle\subseteq K[x_{1},\ldots,x_{n}]$ 是根理想, 当且仅当 $f$ 不被任何非常数多项式的平方整除.
  \item\label{ex:strict-hierarchy-of-ideals}
    \begin{enumerate}[label=(\alph*)]
      \item 举出一个不是根理想的理想 $I\subseteq K[x,y]$.
      \item 举出一个不是素理想的根理想 $I\subseteq K[x,y]$.
      \item 举出一个不是极大理想的素理想 $I\subseteq K[x,y]$.
    \end{enumerate}
  \item\label{ex:square-free-generators-need-not-give-radical} 举例证明: 即使 $f_{1}$ 和 $f_{2}$ 都无平方因子, 理想 $\langle f_{1},f_{2}\rangle\subseteq K[x_{1},\ldots,x_{n}]$ 也不一定是根理想.
  \item\label{ex:zero-divisors-as-radical-ideal} 证明环 $R$ 中所有零因子的集合是根理想.
  \item\label{ex:unit-ideal-iff-radical-unit-ideal} 设 $R$ 是环, $I\subseteq R$ 是理想. 证明 $\sqrt{I}=R$ 当且仅当 $I=R$.
  \item\label{ex:radical-of-intersection} 设 $R$ 是环, $I,J\subseteq R$ 是理想. 证明
    \[
      \sqrt{I\cap J}=\sqrt{I}\cap\sqrt{J}.
    \]
  \item\label{ex:radical-is-intersection-of-primes} 设 $R$ 是环, $I\subseteq R$ 是理想. 证明 $\sqrt{I}$ 是 $R$ 中所有包含 $I$ 的素理想的交.
\end{BookExercises}

\par\smallskip
{\footnotesize
\noindent\textbf{开放获取.} 本章按照
\href{http://creativecommons.org/licenses/by-nc/4.0/}{Creative Commons Attribution-NonCommercial 4.0 International License}
的条款授权. 该许可允许以任何媒介或格式进行非商业性使用, 分享, 改编, 分发和复制, 但须对原作者和来源作适当署名, 提供 Creative Commons 许可的链接, 并说明是否作了修改.

\noindent
\begin{minipage}[t]{.17\linewidth}
  \vspace{0pt}\includegraphics[width=.9\linewidth]{figures/source-cc-by-nc.png}
\end{minipage}\hfill
\begin{minipage}[t]{.81\linewidth}
  \vspace{0pt}本章的图片或其他第三方材料均纳入本章的 Creative Commons 许可, 除非材料的署名行另有注明. 如果材料未纳入本章的 Creative Commons 许可, 且您拟进行的使用未获法律规定的许可或超出允许使用的范围, 则须直接向著作权持有人取得许可.
\end{minipage}\par}
